\documentclass[UTF8]{ctexart}
\usepackage{xeCJK}
\usepackage{geometry}
\geometry{left=2cm,right=2cm,top=2.5cm,bottom=2.5cm}
\setlength{\headheight}{12.7pt}

\usepackage{titlesec}
\titleformat{\section}
  {\normalfont\large\bfseries}
  {\thesection}
  {1em}
  {}
\titlespacing*{\section}
  {0pt}
  {3.5ex plus 1ex minus .2ex}
  {2.3ex plus .2ex}

\usepackage{amsmath}
\usepackage{amssymb}
\usepackage{amsthm}
\usepackage{enumitem}
\setlist[enumerate]{itemsep=0pt,topsep=0pt,parsep=0pt,leftmargin=0pt,itemindent=2.5em}
\setlist[itemize]{itemsep=0pt,topsep=0pt,parsep=0pt,leftmargin=0pt,itemindent=2.5em}
\usepackage{fancyhdr}
\pagestyle{fancy}
\fancyhf{}
\usepackage{graphicx}
\usepackage{hyperref}
\usepackage{indentfirst}
\usepackage{lmodern}


\begin{document}
\setlength{\abovedisplayskip}{1pt}
\setlength{\belowdisplayskip}{1pt}

\section{群的同构定理}

\hypertarget{sec:定理1.1}{}
\noindent\textbf{定理1.1 (群同构第一定理)}

任给群同态$\varphi:G\to G'$, 有群同构
\[
    \overline{\varphi}:\,\!^{\displaystyle G}\!
    /_{\displaystyle\mathrm{Ker}\,\varphi}\cong\mathrm{Im}\,\varphi,
\]
并且对任意的$g\in G$有
$\overline{\varphi}(\,g\cdot\mathrm{Ker}\,\varphi)=\varphi(g)$.

\noindent{\kaishu 证明.}

因为$\mathrm{Ker}\,\varphi\lhd G$,
由\href{04 正规子群与商群#nameddest=sec:定理2.1}{商群的泛性质}, 有唯一的群同态
$\overline{\varphi}:\,\!^{\displaystyle G}\!
/_{\displaystyle\mathrm{Ker}\,\varphi}\to G'$, 并且
\vspace{3pt}
\[
    \mathrm{Ker}\,\overline{\varphi}=\,\!^{\displaystyle\mathrm{Ker}\,\varphi}\!
    /_{\displaystyle\mathrm{Ker}\,\varphi}=\{\mathrm{Ker}\,\varphi\},\quad
    \mathrm{Im}\,\overline{\varphi}=\mathrm{Im}\,\overline{\varphi},
\]
所以$\overline{\varphi}:\,\!^{\displaystyle G}\!
/_{\displaystyle\mathrm{Ker}\,\varphi}\cong\mathrm{Im}\,\varphi$.
\hfill $\qedsymbol$

\vspace{15pt}

\hypertarget{sec:定理1.2}{}
\noindent\textbf{定理1.2 (群同构第二定理)}

任给群$G$, $H<G$和$K\lhd G$, 有$H\cap K\lhd H$, $K\lhd HK$以及群同构
\vspace{3pt}
\[
    \psi:\,\!^{\displaystyle H}\!/_{\displaystyle H\cap K}
    \cong\,\!^{\displaystyle HK}\!/_{\displaystyle K},\\[3pt]
\]
并且对任意的$h\in H$有$\psi(h\cdot(H\cap K))=hK$.

\noindent{\kaishu 证明.}

对任意的$h_1,h_2\in H$和$k_1,k_2\in K$, 有
\[
    (h_1k_1)(h_2k_2)^{-1}=h_1k_1k_2^{-1}h_2^{-1}=
    (h_1h_2^{-1})(h_2k_1k_2^{-1}h_2^{-1})\in H\cdot h_2Kh_2^{-1}=HK,
\]
所以$HK<G$, 从而$K\lhd HK$. 考虑群同态
\[
    \varphi:H\to\,\!^{\displaystyle HK}\!/_{\displaystyle K},
    \quad h\mapsto hK,\\[3pt]
\]
显然$\mathrm{Ker}\,\varphi=H\cap K$, $\mathrm{Im}\,\varphi=
\,\!^{\displaystyle HK}\!/_{\displaystyle K}$. 应用群同构第一定理即得
$\psi:\,\!^{\displaystyle H}\!/_{\displaystyle H\cap K}
\cong\,\!^{\displaystyle HK}\!/_{\displaystyle K}$. \hfill $\qedsymbol$

\vspace{15pt}

\hypertarget{sec:定理1.3}{}
\noindent\textbf{定理1.3 (群同构第三定理)}

任给群$G$, $H\lhd G$, 以及$H,G$的公共正规子群$N$, 有
$\,\!^{\displaystyle H}\!/_{\displaystyle N}\lhd
\,\!^{\displaystyle G}\!/_{\displaystyle N}$以及群同构
\vspace{6pt}
\[
    \psi:\,\!^{\displaystyle G/N}\!/_{\displaystyle H/N}
    \cong\,\!^{\displaystyle G}\!/_{\displaystyle H},\\[3pt]
\]
并且对任意的$g\in G$有$\psi\left(gN\cdot
\,\!^{\displaystyle H}\!/_{\displaystyle N}\right)=gH$.

\noindent{\kaishu 证明.}

首先有投影同态$\pi:G\to\,\!^{\displaystyle G}\!/_{\displaystyle H}$, 满足
\vspace{3pt}
\[
    \mathrm{Ker}\,\pi=H,\quad
    \mathrm{Im}\,\pi=\,\!^{\displaystyle G}\!/_{\displaystyle H},\quad
    \forall g\in G:\pi(g)=gH,\\[3pt]
\]
因为$N<H$, 由商群的泛性质,
有群同态$\varphi:\,\!^{\displaystyle G}\!/_{\displaystyle N}\to
\,\!^{\displaystyle G}\!/_{\displaystyle H}$, 满足
\vspace{6pt}
\[
    \mathrm{Ker}\,\varphi=\,\!^{\displaystyle H}\!/_{\displaystyle N},\quad
    \mathrm{Im}\,\varphi=\,\!^{\displaystyle G}\!/_{\displaystyle H},\quad
    \forall g\in G:\varphi(gN)=\pi(g)=gH,\\[6pt]
\]
再应用群同构第一定理即得群同构$\psi:\,\!^{\displaystyle G/N}\!/_{\displaystyle H/N}
\cong\,\!^{\displaystyle G}\!/_{\displaystyle H}$. \hfill $\qedsymbol$





\newpage

\hypertarget{sec:定理1.4}{}
\noindent\textbf{定理1.4 (子群对应定理)}

任给群$G$和$N\lhd G$, 则有双射
\vspace{3pt}
\[
    \{H<G\mid N<H\}\to\left\{\,\,\!^{\displaystyle G}\!
    /_{\displaystyle N}\text{\,的子群\,}\right\},\quad H\mapsto
    \,\!^{\displaystyle H}\!/_{\displaystyle N},
\]
并且该对应满足:
\begin{enumerate}
    \item 对任意的包含$N$的$G$的子群$H_1,H_2$, $H_1<H_2$当且仅当
    $\,\!^{\displaystyle H_1}\!/_{\displaystyle N}<
    \,\!^{\displaystyle H_2}\!/_{\displaystyle N}$,
    \vspace{3pt}
    \item 对任意的包含$N$的$G$的子群$H_1,H_2$, $H_1\lhd H_2$当且仅当
    $\,\!^{\displaystyle H_1}\!/_{\displaystyle N}\lhd
    \,\!^{\displaystyle H_2}\!/_{\displaystyle N}$,
\end{enumerate}

\noindent{\kaishu 证明.}

\begin{enumerate}
    \item 显然$H_1<H_2$时有
    $\,\!^{\displaystyle H_1}\!/_{\displaystyle N}<
    \,\!^{\displaystyle H_2}\!/_{\displaystyle N}$, 反之当
    $\,\!^{\displaystyle H_1}\!/_{\displaystyle N}<
    \,\!^{\displaystyle H_2}\!/_{\displaystyle N}$时, 对任意的$h_1\in H$有
    \vspace{3pt}
    \[
        h_1N\in\,\!^{\displaystyle H_1}\!/_{\displaystyle N}\subset
        \,\!^{\displaystyle H_2}\!/_{\displaystyle N},
    \]
    从而$h_1\in H_2$, 即$H_1<H_2$.

    \item 当$H_1\lhd H_2$时, 由群同构第三定理得
    $\,\!^{\displaystyle H_1}\!/_{\displaystyle N}\lhd
    \,\!^{\displaystyle H_2}\!/_{\displaystyle N}$,
    
    \vspace{3pt}\hspace{2.17em}
    反之当$\,\!^{\displaystyle H_1}\!/_{\displaystyle N}\lhd
    \,\!^{\displaystyle H_2}\!/_{\displaystyle N}$时,
    对任意的$h_1\in H_1$和$h_2\in H_2$有
    \[
        h_2h_1h_2^{-1}N=h_2N\cdot h_1N\cdot(h_2N)^{-1}\in
        \,\!^{\displaystyle H_1}\!/_{\displaystyle N},
    \]
    从而$h_2h_1h_2^{-1}\in H_1$, 即$H_1\lhd H_2$.

    \vspace{3pt}
    \item 由第一条, 显然$H\mapsto\,\!^{\displaystyle H}\!/_{\displaystyle N}$
    是单射,

    \vspace{3pt}
    \item 对$\,\!^{\displaystyle G}\!/_{\displaystyle N}$的任意子群$K$,
    令$\displaystyle H=\bigcup K$.

    \vspace{3pt}\hspace{2.17em}
    第一, 对任意的$h_1,h_2\in H$, 有$h_1N,h_2N\in K$, 从而
    \[
        h_1h_2^{-1}N=(h_1N)(h_2N)^{-1}\in K\quad\implies\quad h_1h_2^{-1}\in H,
    \]
    即$H<G$.

    \hspace{2.17em}
    第二, 对任意的$n\in N$, 有$nN=N\in K$, 从而$n\in H$. 即$N<H$.

    \hspace{2.17em}
    第三, 对任意的$g\in G$有
    \[
        gN\in\,\!^{\displaystyle H}\!/_{\displaystyle N}\iff
        g\in H\iff g\in\bigcup K\iff gN\in K,\\[3pt]
    \]
    即$\,\!^{\displaystyle H}\!/_{\displaystyle N}=K$, 所以
    $H\mapsto\,\!^{\displaystyle H}\!/_{\displaystyle N}$是满射. \hfill $\qedsymbol$
\end{enumerate}

\end{document}